\begin{korollar}{Gordon'sche Paare und subeuklidische Bildhalbnormen}
{GordonschePaareUndSubeuklidischeBildhalbnormen}
Seien $(E, \norm\cdot_E)$ ein normierter $\R$--Vektorraum, $0\neq x \in E^N$,
$G_x$ der Gauß'sche Operator zum Tupel $(x,g)$ und $\emptyset\neq T
\subseteq S^{m-1}$.\\
Es gelte $\norm\cdot_E \circ \bprod \cdot x_N \le \norm\cdot_{2, N}$.\\
Dann sind nach~\refclaim{IntegrierbarkeitVonExtremaInGordonschenPaaren}
{IntegrierbarkeitVonExtremaInGordonschenPaaren1}
\[\norm\cdot_E \circ \bprod \cdot x _N \circ \gamma, \quad
\inf_{t\in T}\bprod t \delta_m, \quad \sup_{t\in T}\bprod t \delta_m\]
und nach nach~\refclaim{IntegrierbarkeitVonExtremaInGordonschenPaaren}
{IntegrierbarkeitVonExtremaInGordonschenPaaren2}
\[\inf_{t \in T}\, \norm\cdot_E \circ ev(t,\cdot)\circ G_x, \quad
\sup_{t \in T}\, \norm\cdot_E \circ ev(t,\cdot)\circ G_x \]
integrierbar und es gilt:
\begin{enumerate}[(1)]
	\item \label {GordonschePaareUndSubeuklidischeBildhalbnormen1}
\begin{align*}
& \E\left( \norm\cdot_E \circ \bprod \cdot x _N \circ \gamma\right)
	+ \E\left(\inf_{t\in T}\bprod t \delta_m \right)\\
	\le &\,\E\left(\inf_{t \in T} \norm\cdot_E \circ ev(t,\cdot)\circ G_x \right)
	\le \E\left(\sup_{t \in T} \norm\cdot_E \circ ev(t,\cdot)\circ G_x \right)\\
	\le &\,\E\left( \norm\cdot_E \circ \bprod \cdot x _N \circ \gamma\right)
	+ \E\left(\sup_{t\in T}\bprod t \delta_m \right)\text.
\end{align*}
	\item \label {GordonschePaareUndSubeuklidischeBildhalbnormen2}
	Ist $T = - T$, so gilt
	$\E\left(\Inf{t\in T}\bprod t \delta_m \right)
	= - \E\left(\Sup{t\in T}\bprod t \delta_m \right)$.
\end{enumerate}
\end{korollar}
\begin{proof}
(1) Sei $V\coloneqq  \menge{\pfeil \varphi N (x)}{\varphi \in \kabg{E'}01}$.\\
Dann ist $V$ nichtleer und, wie bereits im Beweis von
\refclaim{IntegrierbarkeitVonExtremaInGordonschenPaaren}
{IntegrierbarkeitVonExtremaInGordonschenPaaren1} gesehen, beschränkt.\\
Sei $v \in V$. Dann gibt es ein $\varphi \in \kabg{E'}01$ mit $v=\pfeil\varphi N(x)$.\\
Ferner gilt:
\begin{align*}
\norm v_{2, N}^2 &= \sprod v v_N = \sprod v {\pfeil \varphi N(x)}_N
\overset{\refclaim{EigenschaftenDesLinearkombinators}
{EigenschaftenDesLinearkombinators2}}=
\varphi\left(\bprod v x_N \right) \le \abs{\varphi\left(\bprod v x_N \right)}\\
&\le \opnorm \varphi E \R\norm{\bprod v x_N}_E\overset{\varphi \in \kabg{E'}01}\le
\norm{\bprod v x_N}_E \overset{\text{Vor.}}\le \norm v _{2, N}
\end{align*}
Damit ist $\norm v _{2, N} \le 1$.\\
Daher ist $c\coloneqq 1$ eine obere Schranke von $\menge{\norm w _{2, N}}{w \in V}$.
\hfill $(\ast)$\\
Offensichtlich ist $0 \in V$. Wegen $x \neq 0$ existiert ein $i \in \haken N$
mit $x_i \neq 0$.\\ Nach~\ref{NormUndDualeNorm} existiert ein 
$\alpha \in \kran{E'}01$ mit $\alpha(x_i) = \norm {x_i}_E$.\\
Definiere $v_n \coloneqq  \frac 1 n \pfeil \alpha N (x)$ für alle $n \in \N$.
Dann ist $(v_n)_{n \in \N}$ eine Nullfolge.\\
Für alle $n \in \N$ gilt wegen $\opnorm \alpha E \R= 1$: $v_n \neq 0$ und
nach Definition von $V$ auch $v_n \in V$, also $v_n \in V\setminus \meng 0$
und damit $0$ ein Häufungspunkt von $V$.\\
Schließlich gilt für alle $y\in\R^N$:
\begin{align*}
\sup_{w\in V} \sprod y w _N 
&= \sup_{\tau \in \kabg{E'}01}\sprod y {\pfeil \tau N (x)}_N
\overset{\refclaim{EigenschaftenDesLinearkombinators}
{EigenschaftenDesLinearkombinators2}}= \sup_{\tau \in \kabg{E'}01} 
\tau\left(\bprod y x_N \right)\\& \overset{\ref{NormUndDualeNorm}}=
\norma{\bprod y x_N}_E\text.
\end{align*}
Damit sind wegen $(\ast)$ die Voraussetzungen von~\ref{GordonschePaareUndNormiterierendeMengen} mit $c=1$ erfüllt und die Behauptung ist dann genau die Aussage
\refclaim{GordonschePaareUndNormiterierendeMengen}
{GordonschePaareUndNormiterierendeMengen1}.\\
(2) wurde bereits in~\refclaim{GordonschePaareUndNormiterierendeMengen}
{GordonschePaareUndNormiterierendeMengen2} bewiesen.
\end{proof}